\section{Conclusions}
\label{sec:conclusions}

The four inheritance relations connect rules from different domains and reduce the 256 ECA to 81 classes. The canonical form of the transition graph exactly distinguishes the 88 classes under permutation and mirror with a single space of 6 cells if both the full graph and its quotient by translations are used, but it gives a binary comparison limited to rules of the same domain.

The eigenvalues and the spectral moments only see the attractors: they do not separate rules that differ in their transients, such as 0 and 8, and they stop at 77 classes. The spectral moments with attractor mass keep their fixed length and add the information of the basins: they distinguish 86 of the 88 classes (87 together with the quotient) and separate rules 0 and 8, although with a small distance (0.007 with $n = 12$) because the difference lies in the first steps of a vector that saturates.

Across domains, the attractor mass gives zero distance to the relations that preserve the transition graph and, in the quotient, also to those that change it by a translation; thus, making the neighbors contiguous seems to preserve the behavior (modulo translations, except for non-contiguous neighborhoods on spaces of even size), but the exact comparison is to keep the neighbors at their original positions: in that case the global function is the same and the distance is always zero, as when an end neighbor is removed. State reduction is not identified, because the configurations with the new state dominate the normalization; the plain spectral moments, in contrast, are invariant under it. Moreover, the vector depends on the size of the space, so it is better to compare on the same size. With that precaution, the distance groups the ECA into four families according to the dominant period of their attractors, although the chaotic rules are scattered because their long cycles do not fit in the window.

Future work includes: normalizing the attractor mass so that it is invariant under state reduction, for example by restricting it to the configurations that survive the first step; giving more weight to the first steps and keeping the magnitude of the vector; choosing the space size and the window $K$ systematically; checking whether the similarity holds when the neighbors are made contiguous; studying dimension reduction; and extending the comparison to larger domains through sampling.
